% from ega4-20-fr.tex
\oldpage[IV-4]{249}
tels que $u(a)=x$, $u(b)\in U$ et que la $S$-application rationnelle $\omega\circ u$ soit partout définie.

En effet, par hypothèse toutes les fibres du morphisme projection $X\times_S Y\to X$ sont formées de points isolés, et pour appliquer \sref{IV.20.4.4} il suffit de savoir que la fibre $p^{-1}(x)$ est non vide dans $\Gamma_\omega$.
Or si $h_1$ est l'unique morphisme de la classe $\omega\circ u$, l'unique point de $X\times_S Y$ au-dessus de $x$ et de $h_1(a)$ appartient à $\Gamma_\omega$, d'où la conclusion.
\end{corollary}

\begin{proposition}[20.4.11]
\label{IV.20.4.11}
Soient $X$ un préschéma localement noethérien, $Y$ un $S$-préschéma affine sur $S$, $U$ un ouvert de $X$, $Z=X-U$.
Supposons que l'on ait $\operatorname{prof}_Z(\sh{O}_X)\geq 2$ \sref{IV.5.10.1} ; alors tout $S$-morphisme $f:U\to Y$ se prolonge de façon unique en un $S$-morphisme de $X$ dans $Y$.
\end{proposition}

\begin{proof}
On peut se borner au cas où $S$ et $X$ sont affines et (en vertu de \sref[I]{I.3.3.14}) au cas où $S=X$ ; on a donc $X=\Spec(A)$, $Y=\Spec(B)$, $B$ étant une $A$-algèbre de type fini.
Comme $B$ est quotient d'une algèbre de polynômes $A[(T_\lambda)]_{\lambda\in L}$, $Y$ est un sous-préschéma fermé de $Y'=X[(T_\lambda)]_{\lambda\in L}$.
D'autre part, l'hypothèse sur $Z$ entraîne que $U$ est schématiquement dense dans $X$ en vertu de \sref{IV.20.2.13}[(iv)] et \sref{IV.5.10.2}.
Si l'on prouve que tout $X$-morphisme $u:U\to Y$ se prolonge de façon unique en un $X$-morphisme $v':X\to Y'$, il en résultera que $v'$ se factorisera en $X\xrightarrow{v}Y\to Y'$ ; en effet, le sous-préschéma $v'^{-1}(Y)$ est fermé et majore $U$ \sref[I]{I.4.4.1}, donc est identique à $X$ \sref{IV.20.3.8.8}.
Dans ces conditions, $v$ sera l'unique $S$-morphisme de $X$ dans $Y$ prolongeant $u$.
On peut donc se borner au cas où $Y=Y'$.
Mais alors il y a correspondance biunivoque entre les $X$-morphismes d'un ouvert $U\subset X$ dans $Y'$ et les familles $(s_\lambda)_{\lambda\in L}$ de sections de $\sh{O}_X$ au-dessus de $U$ \sref[II]{II.1.7.9} ; la conclusion résulte donc de \sref{IV.5.10.5}.
\end{proof}

\begin{corollary}[20.4.12]
\label{IV.20.4.12}
Soient $X$ un $S$-préschéma localement noethérien réduit et vérifiant la condition $(S_2)$ \sref{IV.5.7.2} (par exemple un $S$-préschéma localement noethérien et normal \sref{IV.5.8.6}), $Y$ un $S$-préschéma affine sur $S$, $f$ une $S$-application rationnelle de $X$ dans $Y$ ; alors toute composante irréductible de $X-\operatorname{dom}(f)$ est de codimension $1$ dans $X$.
\end{corollary}

\begin{proof}
Il revient au même de dire que si $Z^{(2)}$ est l'ensemble des $x\in X$ tels que $\dim(\sh{O}_{X,x})\geq 2$, alors, pour toute partie fermée $Z\subset Z^{(2)}$ de $X$, tout $S$-morphisme de $X-Z$ dans $Y$ se prolonge en un $S$-morphisme de $X$ dans $Y$ ; or, cela résulte de l'hypothèse sur $X$ \sref{IV.5.7.2} et de \sref{IV.20.4.11}.
\end{proof}

\subsection{Pseudo-morphismes relatifs}
\EGAReaderTarget{ega4.c04.subsec.20-5}{IV.20.5}

\begin{env}[20.5.1]
\EGAReaderTarget{ega4.c04.statement.20-5-1}{IV.20.5.1}
Soient $X$, $Y$ deux $S$-préschémas.
Il résulte des définitions \sref{IV.11.10.8} que l'intersection de deux ouverts $U$, $U'$ de $X$, \emph{universellement schématiquement denses par rapport à $S$}, possède encore cette propriété.
On définit donc une relation d'équivalence entre $S$-morphismes $u:U\to Y$ en remplaçant dans \sref{IV.20.2.1} « schématiquement dense » par « universellement schématiquement dense par rapport à $S$ ».
Une classe d'équivalence suivant cette relation s'appelle un \emph{pseudo-morphisme de $X$ dans $Y$ relativement à $S$}, et l'ensemble de ces classes se note $\operatorname{Ps.hom}_{X/S}(X,Y)$.
\end{env}

\begin{env}[20.5.2]
\label{IV.20.5.2}
Comme tout ouvert de $X$ universellement schématiquement dense par rapport à $S$ est en particulier schématiquement dense, les éléments d'un pseudo-morphisme
